\scanpage{38}
\mspage{\handwritten{100}}

Faut-il croire en l'homéopathie vétérinaire ?

Un dossier de Jacinte Giscard

Sur les 5.000 vétérinaires exerçant en France\unsure{.|,} 4.000
\unsure{utilisent} en certaines occasions des préparations
\unsure{homéopathiques} et 400 d'entre eux traitent 30 \% de
\unsure{leurs} patients uniquement par cette pratique. Les chiffres le révèlent : l'homéopathie vétérinaire est une réalité bien palpable.

Pourtant\unsure{.|,} elle est loin de faire l'unanimité. C'est
même bien souvent un objet de polémique entre partisans qui la portent aux nues et détracteurs qui la font
descendre aux enfers. Parce qu'elle ne s'assoie pas
encore suffisamment sur des preuves scientifiques
indiscutables, l'homéopathie est encore une médecine
qui dérange. Cela est rendu plus vrai du fait qu'elle
s'entoure parfois d'un parfum d'ésotérisme un peu trop
discréditant.

Vous êtes nombreux à nous interroger dans votre
courrier sur l'homéopathie et sur son utilisation chez
l'animal. ``Où en est l'homéopathie vétérinaire, quels
développements peut-on en attendre à l'avenir ?'' sont
des questions qui reviennent chaque jour.

Pour vous, nous avons donc réalisé ce grand dossier, un peu technique parfois (c'était inévitable), mais
riche d'enseignements. Il ne faut pas croire en l'homéopathie comme en une panacée universelle. Mais il est
indubitable qu'elle donne, en de nombreuses occasions, de bons résultats. Essayons de comprendre
pourquoi.

\mspage{47}

\handwritten{``30 millions d'amis''}

\handwritten{nov. \unsure{88|85}}

En fait, et il faut bien se le
mettre dans la tête, l'homéopathie vétérinaire a un intérêt dans toutes les formes
d'intervention médicale. Elle
peut s'appliquer efficacement tant aux cas d'urgence
comme les diarrhées suraiguës ou les traumatismes,
qu'aux maladies courantes
et à la prévention sanitaire.

% archive sheet number: 37
